% 本文件由 LaTeX 排版 Agent 根据有效 translation.json 自动生成。
% author=Peter of Alexandria; work_id=0621; title=Fragments

\chapter{\WorkTitle{Fragments}}

% structure: p0001 source=h1 target=omitted reason=page-title-already-rendered-by-parent

% structure: p0002 source=h2 target=section reason=relative-heading-rank; base=h2

\section{一、致亚历山大里亚教会书}

伯多禄，致在主内蒙爱、立于天主信德中的弟兄们，愿主赐予平安。\newline{}因我已得知梅里提乌行事全然不为公益——他既不满足于至圣主教及殉道者们的信函——反而侵入我的教区，擅自僭权，试图使司铎们及受托探访贫苦者的人脱离我的权下；他并显出其好为首之心，在狱中私自为自己祝圣了几人。\newline{}因此，请留意此事，勿与他共融，直到我偕同几位明智谨慎之士与他相会，察看他所图谋的计策为何。\newline{}再见。\par

% structure: p0004 source=h2 target=section reason=relative-heading-rank; base=h2

\section{二、论天主性}

因为确实，\BibleQuote{恩宠和真理是由耶稣基督而来的}，\BibleRef{若 1:17} 因此我们也藉恩宠得救，正如宗徒所说：\BibleQuote{也不是出于自己，也不是出于行为，免得任何人自夸}，\BibleRef{弗 2:8} 按天主的旨意，\BibleQuote{圣言成了血肉}，且\BibleQuote{以人的形态出现}。\BibleRef{斐 2:7} 但祂并未失去祂的天主性。因为\BibleQuote{祂本是富有的，却成了贫穷的}，\BibleRef{格后 8:9} 并非要完全与祂的能力和光荣分离，而是为使我们这些罪人，祂亲自承担死亡，义者为不义者，为将我们带到天主面前，\BibleQuote{在肉身内被处死，却在圣神内复活}；其后还有别的话。因此福音作者也证实真理说：\BibleQuote{圣言成了血肉，居住在我们中间}；那时，自天使问候童贞女说：\BibleQuote{万福！充满恩宠者，主与你同在}。当加俾额尔说\BibleQuote{主与你同在}时，意指天主圣言与你同在。因为他显示祂在母胎中受孕，并将成为血肉；正如经上所记：\BibleQuote{圣神要临于你，至高者的能力要荫庇你，因此，那要诞生的圣者，将称为天主子}；\BibleRef{路 1:35} 其后还有别的话。天主圣言，在没有男人参与的情况下，按天主的旨意——祂轻易成就一切——在童贞女的胎中成为血肉，不需要有男人的临在与行动。因为比男人更有效的是天主的能力荫庇童贞女，连同临于她的圣神。\par

% structure: p0006 source=h2 target=section reason=relative-heading-rank; base=h2

\section{三、论我们救主的来临}

祂对犹达斯说：\BibleQuote{你用亲吻出卖天主子吗？}\BibleRef{路 22:48} 这些事以及类似的事，还有祂所显的一切征兆和奇迹，都证明祂是成了人的天主。因此两件事得以显明：祂按本性是天主，祂按本性也是人。\par

% structure: p0008 source=h2 target=section reason=relative-heading-rank; base=h2

\section{四、论基督与我们同住}

因此两件事得以证明：祂按本性是天主，并按本性成了人。\par

% structure: p0010 source=h2 target=section reason=relative-heading-rank; base=h2

\section{五、论直到耶路撒冷毁灭时，犹太人正确规定第一个太阴月的第十四日}

% structure: p0011 source=h3 target=subsection reason=relative-heading-rank; base=h2

\subsection{1}

1. 既然天主的仁慈处处伟大，我们当赞美祂，也因祂赐给我们真理之神，引导我们进入一切真理。为此，法律规定了阿彼布月为正月，并为我们显明它为年中的第一月；无论是住在那以前的古代作者，还是耶路撒冷毁灭之后的后期作者，都指明它拥有最清晰且显然确定的时期，尤其是在某些地方收割较早，有时较晚，以至于有时在逾越节之前，有时在之后，正如在最初颁布法律、逾越节之前所发生的那样，按经上所记：\BibleQuote{但是小麦和粗麦却没有被打坏，因为还没有长成。}\BibleRef{出 9:32} 因此法律正确地规定，从春分起，第一月的第十四日落在哪一周，逾越节就应在那一周庆祝，并先以相称合宜的赞美诗歌来庆祝。因为祂说：\BibleQuote{这月必为你们之正月}，\BibleRef{出 12:2} 此时太阳在夏季发出远为强烈而清晰的光，白昼延长变长，黑夜缩短变短。再者，当新苗长出时，它们被完全净除，运到打谷场；不仅如此，所有灌木开花绽放。因此立刻发现它们交替结出各种不同的果实，以致那时能找到葡萄串；正如立法者所说：\BibleQuote{那时正是春天，是葡萄初熟的季节}；\BibleRef{户 12:24} 当他派人窥探那地时，他们用杠抬回一大串葡萄，还有石榴和无花果。因为那时，如他们所说，我们永恒的天主，万有的创造者与造物主，创造了万物，并对它们说：\BibleQuote{愿地生出青草、结种子的菜蔬、结果实的果树，各按其类，果实里有种子，在地上}。然后祂加上：\BibleQuote{事就这样成了；天主看了认为好。}此外，祂清楚显明，希伯来人中第一月是由法律规定的，我们知道犹太人遵守它直到耶路撒冷毁灭，因为这是希伯来传统所传下来的。但城被毁后，它受到某些人心硬者的讥诮，我们却按法律真诚地领受了它；并在此中，按圣言，当祂谈及我们神圣庆节的日子时，那蒙拣选者得到了它；其余的人却心硬了，正如经上所说；之后还有别的话。\par

2. 祂接着这样说：\Quote{凡他们为我的名对你们所行的这一切，都是因为他们不认识那派遣我来的。} \BibleRef{若 15:21} 但是，若他们不认识那派遣者与被派遣者，则无可怀疑，他们对于法律所规定的逾越节是无知的，以致不仅在选择地方上犯了错误，而且在计算每月之首（即一年中诸月之首）上也犯了错误；在这一月的第十四日，经准确观察后，在春分之后，古人依照天主的命令庆祝逾越节；而今人却在春分之前庆祝它，完全是由于疏忽和错误，他们不知道如何在恰当的季节庆祝它，正如那位在这些事上所描述者所承认的那样。\par

3. 因此，犹太人有时按月在法默诺特月\OriginalTerm{法默诺特月}{Phamenoth}的月相错误地庆祝他们的逾越节，或按每三年一次的闰月，在法尔穆提月\OriginalTerm{法尔穆提月}{Pharmuthi}庆祝，这对我们无关紧要。因为我们唯一的目的就是纪念祂的受难，而且就在此时；正如那些亲眼目睹者从起初，在埃及人信从之前，所传下来的。因为他们不是通过观察月相而必然在法默诺特月十六日庆祝，而是每三年一次在法尔穆提月庆祝；因为从起初，在基督来临之前，他们似乎就这样做了。因此，当主通过先知责备他们时，祂说：\Quote{他们心中常常错误；我就在愤怒中起誓说：他们不得进入我的安息。}\par

4. 因此，如你所见，在这方面，你似乎大大地撒谎，不仅是针对人，也是针对天主。首先，事实上，在这件事上，犹太人从未犯过错误，因为他们与那些亲眼目睹者和服务者交往，更不用说从起初、在基督来临之前了。因为天主并没有说，他们如同你所写的那样，关于逾越节的法律诫命，心中常常错误，而是由于他们所有其他的悖逆，以及由于他们邪恶和不当的行为，当时祂看到他们转向偶像崇拜和奸淫。\par

5. 此外，在几句话之后。因此，即使在这点上，你既已昏睡，就要用训道者的鞭策大大地、非常地唤醒自己，尤其要记起他所说的话：\Quote{在铺石路上滑倒，以及用舌头滑倒。} \BibleRef{德 20:18} 因为，你再次看到，你加在他们领袖身上的指控反弹了回来；不仅如此，人们还可以怀疑随后有极大的危险，因为我们听说，人向高处抛起的石头会落回他自己的头上。在这方面，一个胆敢指控梅瑟（天主的仆人）或继承他的农的儿子若苏厄，或那些依次正确追随他们并统治的人——我是指民长、出现的君王、或圣神所感动的先知、以及那些在大司祭中无可指责的、并跟随传统而未改变任何事、而在其季节遵守逾越节及其它节期的人——的人，更是鲁莽。\par

6. 另外，在其它事情之后。但你本应走一条更稳妥、更吉祥的道路，而不应轻率和诽谤地断言，他们似乎从起初和一直都对逾越节有错误，这是你无法证明的，无论你可能想对当前那些因严重偏离而犯错的人提出什么指控，他们已偏离了关于逾越节及其它事物的法律诫命。因为古人似乎是在春分之后守逾越节的，你若阅读古代著作，尤其是那些由博学的希伯来人所写的，就能发现这一点。\par

7. 因此，直到主受难时期，以及发生在罗马皇帝韦斯巴芗治下的耶路撒冷最后毁灭之时，以色列子民正确遵守第一个太阴月的第十四天，并在这天庆祝法律所规定的逾越节，这一点已简要证明。因此，当圣先知们，以及一切（如我所说）公义正直地行走在主法律中的人，连同全体百姓，庆祝一个预表性和影儿般的逾越节时，一切有形与无形受造物的创造者与主，独生子，与圣父及圣神同永恒的圣言，并与祂们具有同一性体，按祂的神性，我们的主和天主耶稣基督，在世界的末期，按肉体由我们神圣荣耀的主母、天主之母及终身童贞者，确实由天主之母玛利亚所生；并在地上显现，按人性与同祂一性体的人真正实在交往，祂自己也在祂公开传教之前的年和公开传教期间，与百姓一同庆祝了法律所规定的影儿般的逾越节，吃了预表性的羔羊。因为救主自己在福音中说：\Quote{我来不是为废除法律或先知，而是为成全。}\par

但祂在公开传教之后并没有吃逾越节羔羊，而是亲自作为真羔羊在逾越节盛宴中受苦，正如神圣的福音作者\Person{若望}在祂所写的福音中所教导我们的，其中祂这样说：\Quote{然后他们把耶稣从\Person{盖法}那里带到总督府；那时是清晨；他们自己却不进总督府，免得被玷污，但为了能吃逾越节。}\BibleRef{若 18:28}之后再略述。\Quote{\Person{比拉多}听了这话，就带耶稣出来，到了一个地方，名叫铺石地，希伯来语叫\Place{加巴达}，就在审判台上坐下。那天是逾越节的预备日，约在第三时辰}正如正确的抄本所记载的，并且福音作者亲手所写的副本本身，因天主的恩宠，保存在\Place{厄弗所}最神圣的教会中，在那里为信徒所敬拜。同一位福音作者又说：\Quote{犹太人因那日是预备日，免得尸体在安息日留在十字架上（因为那安息日是大日子），就求\Person{比拉多}叫人打断他们的腿，把他们拿去。}\BibleRef{若 19:31}因此，在那一天，犹太人在晚上要吃逾越节，我们的主和救主耶稣基督被钉在十字架上，成为那些要因信德分享关于祂的奥迹之人的祭品，正如真福\Person{保禄}所写的：\Quote{因为我们的逾越节羔羊基督已被祭献了。}\BibleRef{格前 5:7}而不是像某些人，由于无知而冲动，自信地断言祂吃了逾越节后才被出卖；我们既未从圣福音作者学知，也没有任何一位蒙福的宗徒传递给我们。所以，在我们的主和天主耶稣基督按肉身为我们受苦之时，祂并没有吃律法下的逾越节；而是，正如我所说，祂自己作为真羔羊，在预像的逾越节盛宴中为我们被祭献，在预备日，即第一个太阴月的第十四天。因此，预像的逾越节在那时停止了，因为真逾越节已经来到：\Quote{因为我们的逾越节羔羊基督已被祭献了，}如前所述，并且如那蒙选的器皿、宗徒\Person{保禄}所教导的。\par

% structure: p0020 source=h3 target=subsection reason=relative-heading-rank; base=h2

\subsection{2}

现在是预备日，约在第三时辰，正如准确的抄本以及福音作者\Person{若望}亲手所写的副本本身所记载的，该副本因天主的恩宠，至今仍保存在\Place{厄弗所}最神圣的教会中，并为信徒所敬拜。\par

% structure: p0022 source=h2 target=section reason=relative-heading-rank; base=h2

\section{六、论灵魂与身体}

关于从天上来的第二人，就祂的神性和人性而言，我们在以上根据蒙福的宗徒所写的内容已经解释过了；现在我们觉得有必要解释关于第一人，即出于地、属于土的事，即要证明这一点：他是在同一时间被创造的，是同一个，尽管有时他被分别称为外在的人和内在的人。因为，根据救恩之言，那位创造外在的，也创造了内在的，祂确实在同一行动中、在同一时间创造了二者，正是在那一天，天主说：\Quote{让我们按我们的肖像，按我们的模样造人。}\BibleRef{创 1:26}由此显明，人并非由身体与某种预先存在的原型结合而成。因为，如果大地遵循造物主的命令产生了其他有生命的动物，那么天主从地上所取的尘土，更因天主的旨意和行动而接受了生命的能力。\par

% structure: p0024 source=h2 target=section reason=relative-heading-rank; base=h2

\section{七、残篇}

我真不幸！我竟不记得天主洞察心思，聆听灵魂的声音。我明知是罪却转向它，对自己说：天主是仁慈的，祂必容忍我；当我未立刻受打击时，我不但不停止，反而藐视祂的宽容，耗尽了天主的忍耐。\par

% structure: p0026 source=h2 target=section reason=relative-heading-rank; base=h2

\section{八、论圣玛窦}

在《玛窦福音》中，主对出卖祂的人说：\Quote{你用亲吻来出卖人子吗？}殉道者、亚历山大里亚总主教\Person{伯多禄}注释说，这以及其他类似的事：\Quote{祂所显示的一切标记和所行的奇迹，都证明祂是降生成人的天主；因此两件事同时被证明：祂按本性是天主，又按本性成了人。}\par

% structure: p0028 source=h2 target=section reason=relative-heading-rank; base=h2

\section{九、来自一篇讲道}

与此同时，福音作者坚定地说：\Quote{圣言成了血肉，居住在我们中间。}\BibleRef{若 1:14}由此我们得知，当天使以这些话问候童贞玛利亚时：\Quote{万福！充满恩宠者，主与你同在，}\BibleRef{路 1:28}他意在表明天主圣言与你同在，并显示祂要从她的胎中出来，成为血肉，正如所记载的：\Quote{圣神要临于你，至高者的能力要庇荫你，因此那要诞生的圣者，将称为天主子。}\BibleRef{路 1:35}\par

\SourceRecord{Translated by James B.H. Hawkins.；Ante-Nicene Fathers；Vol. 6.；Edited by Alexander Roberts, James Donaldson, and A. Cleveland Coxe.；Buffalo, NY: Christian Literature Publishing Co.,；1886.；<http://www.newadvent.org/fathers/0621.htm>.}
